\documentclass{techbrief}

\title{Entropy is not continuous in standard probability spaces}
\author{Gabriele Carcassi}
\category{Ensemble spaces}
\tags{Entropy, total variation, $L^1$ topology, Jensen-Shannon divergence, probability measures}

\begin{document}
	
	\maketitle
	
	\begin{abstract}
		We show that the topologies coming from $L^1$, total variation and Jensen--Shannon divergence are all equivalent but do not guarantee the continuity of the entropy. However, a topology that guarantees that mixing (i.e. convex combination) and the entropy are continuous must contain the $L^1$/total variation/JSD topology.
	\end{abstract}
	
	\section{Introduction}
	
	Unlike in the finite-dimensional case, multiple topologies are possible in infinite-dimensional linear/convex spaces. This raises the question as to what is the appropriate one for probability spaces, particularly in the context of ensembles with an entropy defined. The $L^1$ topology for probability densities, the total variation topology for probability measures and the topology generated by open balls of the Jensen--Shannon divergence can be shown to be all equivalent, but they are not enough to guarantee the continuity of the entropy. That is, a sequence of distributions may converge in total variation while its entropy fails to converge. Conversely, we show that a topology that guarantees the continuity of statistical mixing and the entropy guarantees the convergence in $L^1$/total variation/JSD.
	
	\section{Definitions and basic properties}
	
	Let $(X,\Sigma,\mu)$ be a measure space. Let $S(\rho) = -\int_X \rho \log \rho\,d\mu$ be the entropy of a probability measure (i.e. $\rho$ is the Radon--Nikodym derivative of a probability measure absolutely continuous with respect to $\mu$). Let $P_\mu = \left\{
	\rho \in L^1(\mu)
	\;\middle|\;
	\rho \geq 0,\;
	\int_X \rho\,d\mu = 1,\;
	S(\rho) \in \mathbb{R}
	\right\}$ be the space of $\mu$-densities with finite entropy.\footnote{Logarithms are taken in base $2$ unless otherwise stated.}
	
	We now want to compare three different convergence criteria that give rise to the same topology. Considering the probability densities as functions, the $L^1$ distance between densities is
	\begin{equation}
		d_1(\rho,\sigma)=\|\rho-\sigma\|_{L^1(\mu)}
		=
		\int_X |\rho-\sigma|\,d\mu.
	\end{equation}
	If $p_\rho$ and $p_\sigma$ are the corresponding probability measures, then the total variation between the probability measures is given by
	\begin{equation}
		\|p_\rho-p_\sigma\|_{\mathrm{TV}}
		=
		\frac{1}{2}\|\rho-\sigma\|_{L^1(\mu)}.
	\end{equation}
	Thus $L^1$ convergence of densities is equivalent to convergence in total variation of the corresponding measures. The Jensen--Shannon divergence between the two densities is given by
	\begin{equation}
		\operatorname{JSD}(\rho,\sigma)
		=
		S\left(\frac{\rho+\sigma}{2}\right)
		-
		\frac{1}{2}S(\rho)
		-
		\frac{1}{2}S(\sigma).
	\end{equation}
	We can show that this too generates the same topology.

\begin{lemma}
	For $u\in[-1,1]$, define
	\begin{equation}
		f(u)
		=
		\frac{1}{2}
		\left[
		(1+u)\log(1+u)
		+
		(1-u)\log(1-u)
		\right],
	\end{equation}
	with the convention $0\log 0=0$. Then
	\begin{equation}
		\frac{u^2}{2\ln 2}
		\leq
		f(u)
		\leq
		|u|.
	\end{equation}
\end{lemma}

\begin{proof}
	The function $f$ is even, with $f(0)=0$ and $f(1)=1$. For $u\in(-1,1)$,
	\begin{equation}
		f''(u)=\frac{1}{(1-u^2)\ln 2}.
	\end{equation}
	Therefore $f''(u)\geq \frac{1}{\ln 2}$. Since $f(0)=f'(0)=0$, this gives
	\begin{equation}
		f(u)\geq \frac{u^2}{2\ln 2}.
	\end{equation}
	Moreover, $f''(u)\geq0$, so $f$ is convex. Hence on $[0,1]$, $f$ lies below the chord joining $(0,0)$ and $(1,1)$:
	\begin{equation}
		f(u) = f( (1-u) 0 + u 1 )\leq (1-u)f(0) + u f(1) = u.
	\end{equation}
	By evenness, this gives $f(u)\leq |u|$ on all of $[-1,1]$.
\end{proof}

\begin{prop}[Jensen--Shannon and total variation generate the same topology]\label{JSDisL1}
	Let $\rho,\sigma\in P_\mu$. Then
	\begin{equation}
		\frac{1}{8\ln 2}\|\rho-\sigma\|_{L^1(\mu)}^2
		\leq
		\operatorname{JSD}(\rho,\sigma)
		\leq
		\frac{1}{2}\|\rho-\sigma\|_{L^1(\mu)}.
	\end{equation}
	Consequently,
	\begin{equation}
		\operatorname{JSD}(\rho_n,\rho)\to0
		\quad\Longleftrightarrow\quad
		\|\rho_n-\rho\|_{L^1(\mu)}\to0.
	\end{equation}
	Equivalently, the topology generated by the Jensen--Shannon divergence is the total variation topology on the associated probability measures.
\end{prop}

\begin{proof}
	Let
	\begin{equation}
		m=\frac{\rho+\sigma}{2}.
	\end{equation}
	Where $\rho+\sigma>0$, define
	\begin{equation}
		u=\frac{\rho-\sigma}{\rho+\sigma},
	\end{equation}
	and set $u=0$ where $\rho+\sigma=0$. Then $u\in[-1,1]$ and
	\begin{equation}
		\rho=m(1+u),
		\qquad
		\sigma=m(1-u).
	\end{equation}
	Substituting into the definition of Jensen--Shannon divergence gives
	\begin{equation}
		\operatorname{JSD}(\rho,\sigma)
		=
		\int_X m f(u)\,d\mu.
	\end{equation}
	By the scalar bounds on $f$,
	\begin{equation}
		\frac{1}{2\ln 2}\int_X m u^2\,d\mu
		\leq
		\operatorname{JSD}(\rho,\sigma)
		\leq
		\int_X m|u|\,d\mu.
	\end{equation}
	Since $m\,d\mu$ is a probability measure, Cauchy--Schwarz gives
	\begin{equation}
		\begin{aligned}
			\left(\int_X m|u|\,d\mu\right)^2 &= \left(\int_X \sqrt{m} \left(\sqrt{m}|u|\right)\,d\mu\right)^2
			\leq \left(\int_X \sqrt{m}^2 \,d\mu\right)
			\left(\int_X \left(\sqrt{m}|u|\right)^2\,d\mu\right) \\
			&=\left(\int_X m\,d\mu\right)
			\left(\int_X m u^2\,d\mu\right)
			=
			\int_X m u^2\,d\mu,
		\end{aligned}
	\end{equation}
	Moreover,
	\begin{equation}
		\int_X m|u|\,d\mu
		=
		\int_X \frac{\rho+\sigma}{2}
		\frac{|\rho-\sigma|}{\rho+\sigma}\,d\mu
		=
		\frac{1}{2}\|\rho-\sigma\|_{L^1(\mu)}.
	\end{equation}
	Putting it all together, we have
	\begin{equation}
		\begin{aligned}
		\frac{1}{2\ln 2}\left(\int_X m|u|\,d\mu\right)^2
&\leq
\operatorname{JSD}(\rho,\sigma)
\leq
\int_X m|u|\,d\mu \\
		\frac{1}{8\ln 2}\|\rho-\sigma\|_{L^1(\mu)}^2
&\leq
\operatorname{JSD}(\rho,\sigma)
\leq
\frac{1}{2}\|\rho-\sigma\|_{L^1(\mu)}. \\
		\end{aligned}
	\end{equation}
	The two inequalities imply the equivalence of Jensen--Shannon convergence and $L^1$ convergence. Since $L^1$ convergence of densities is equivalent to total variation convergence of the associated probability measures, the Jensen--Shannon topology is the total variation topology.
\end{proof}

We want to stress that, given the measure-theoretic construction, all these results apply uniformly to both countable discrete spaces with counting measure and continuous dominated measure spaces.
	
\section{Failure of entropy continuity}

We now show that, whenever the measure space contains finite-measure sets of arbitrarily large measure, entropy is not continuous with respect to the $L^1$ topology.

\begin{thrm}[Entropy is not $L^1$-continuous on spaces with arbitrarily large finite-measure sets]
	Let $(X,\Sigma,\mu)$ be a measure space. Suppose there exists a measurable set $B\in\Sigma$ such that
	\begin{equation}
		0<\mu(B)<\infty,
	\end{equation}
	and a sequence of measurable sets $A_n\in\Sigma$ such that
	\begin{equation}
		A_n\cap B=\emptyset,\qquad
		1<\mu(A_n)<\infty,\qquad
		\mu(A_n)\to\infty.
	\end{equation}
	Then the entropy
	\begin{equation}
		S(\rho)=-\int_X \rho\log\rho\,d\mu
	\end{equation}
	is not continuous in the $L^1(\mu)$ topology on finite-entropy probability densities.
\end{thrm}

\begin{proof}
	For a measurable set $U$ with $0<\mu(U)<\infty$, define the uniform probability density on $U$ by
	\begin{equation}
		u_U=\frac{\mathbf{1}_U}{\mu(U)}.
	\end{equation}
	Its entropy is
	\begin{equation}
		S(u_U)
		=
		-\int_U \frac{1}{\mu(U)}
		\log\left(\frac{1}{\mu(U)}\right)d\mu
		=
		\log \mu(U).
	\end{equation}
	
	Let
	\begin{equation}
		\rho=u_B.
	\end{equation}
	Then
	\begin{equation}
		S(\rho)=\log\mu(B).
	\end{equation}
	
	Define
	\begin{equation}
		\varepsilon_n=\frac{1}{\log\mu(A_n)}.
	\end{equation}
	Since $\mu(A_n)\to\infty$, we have $\varepsilon_n\to0$. Passing to a tail of the sequence if necessary, we may assume $\varepsilon_n\in(0,1)$ for all $n$.
	
	Now define
	\begin{equation}
		\rho_n
		=
		(1-\varepsilon_n)u_B
		+
		\varepsilon_n u_{A_n}.
	\end{equation}
	Since $A_n$ and $B$ are disjoint,
	\begin{equation}
		\|\rho_n-\rho\|_{L^1}
		=
		\varepsilon_n\|u_{A_n}\|_{L^1}
		+
		\varepsilon_n\|u_B\|_{L^1}
		=
		2\varepsilon_n
		\to0.
	\end{equation}
	Thus $\rho_n\to\rho$ in $L^1(\mu)$.
	
	On the other hand, since the supports are disjoint, entropy satisfies the exact mixture formula
	\begin{equation}
		S(\rho_n)
		=
		h_2(\varepsilon_n)
		+
		(1-\varepsilon_n)S(u_B)
		+
		\varepsilon_n S(u_{A_n}),
	\end{equation}
	where
	\begin{equation}
		h_2(\varepsilon)
		=
		-\varepsilon\log\varepsilon
		-
		(1-\varepsilon)\log(1-\varepsilon).
	\end{equation}
	Therefore
	\begin{equation}
		S(\rho_n)
		=
		h_2(\varepsilon_n)
		+
		(1-\varepsilon_n)\log\mu(B)
		+
		\varepsilon_n\log\mu(A_n).
	\end{equation}
	By construction,
	\begin{equation}
		\varepsilon_n\log\mu(A_n)=1.
	\end{equation}
	Moreover,
	\begin{equation}
		h_2(\varepsilon_n)\to0
	\end{equation}
	and
	\begin{equation}
		(1-\varepsilon_n)\log\mu(B)\to\log\mu(B).
	\end{equation}
	Hence
	\begin{equation}
		S(\rho_n)\to\log\mu(B)+1.
	\end{equation}
	But
	\begin{equation}
		S(\rho)=\log\mu(B).
	\end{equation}
	Thus $\rho_n\to\rho$ in $L^1(\mu)$, while $S(\rho_n)\not\to S(\rho)$. Hence entropy is not $L^1$-continuous.
\end{proof}

\begin{remark}
	The proof only uses the existence of finite-measure sets of arbitrarily large measure, disjoint from a fixed finite-measure set. In the countable discrete case with counting measure, one may take $B=\{1\}$ and $A_n$ to be finite sets of cardinality tending to infinity, disjoint from $B$. Then the construction reduces to the standard example where a probability mass $\varepsilon_n$ is spread uniformly over many points.
\end{remark}

\begin{remark}
	Choosing instead
	\begin{equation}
		\varepsilon_n=\frac{1}{\sqrt{\log\mu(A_n)}},
	\end{equation}
	we still have
	\begin{equation}
		\varepsilon_n\to0,
	\end{equation}
	and therefore
	\begin{equation}
		\rho_n\to\rho
		\qquad\text{in }L^1(\mu).
	\end{equation}
	However,
	\begin{equation}
		\varepsilon_n\log\mu(A_n)
		=
		\sqrt{\log\mu(A_n)}
		\to\infty.
	\end{equation}
	Therefore
	\begin{equation}
		S(\rho_n)\to\infty
	\end{equation}
	while still
	\begin{equation}
		\rho_n\to\rho
		\qquad\text{in }L^1(\mu).
	\end{equation}
	Thus arbitrarily small $L^1$ perturbations can carry arbitrarily large entropy.
\end{remark}

\section{Entropy and mixing force total variation}

We now show that requiring entropy and statistical mixing to be continuous forces the topology to contain the $L^1$/total variation/JSD topology.

\begin{thrm}[Entropy continuity and mixing continuity force total variation]
	Let $\tau$ be a topology on $P_\mu$ such that:
	\begin{itemize}
		\item statistical mixing (e.g. $(\lambda,\rho,\sigma)
		\mapsto
		\lambda\rho+(1-\lambda)\sigma$) is continuous;
		\item entropy (i.e. $S(\rho)=-\int_X \rho\log\rho\,d\mu$) is continuous.
	\end{itemize}
	Then $\tau$ contains the topology generated by the Jensen--Shannon divergence. Consequently, $\tau$ contains the total variation topology, equivalently the $L^1(\mu)$ topology on densities.
\end{thrm}

\begin{proof}
	For $\rho,\sigma\in P_\mu$, the Jensen--Shannon divergence is
	\begin{equation}
		\operatorname{JSD}(\rho,\sigma)
		=
		S\left(\frac{\rho+\sigma}{2}\right)
		-
		\frac{1}{2}S(\rho)
		-
		\frac{1}{2}S(\sigma).
	\end{equation}
	Since statistical mixing is continuous and entropy is continuous, the map
	\begin{equation}
		(\rho,\sigma)
		\mapsto
		\operatorname{JSD}(\rho,\sigma)
	\end{equation}
	is continuous. Therefore $\tau$ must contain the topology generated by the Jensen--Shannon divergence which, as shown in Prop. \ref{JSDisL1}, is the same as the total variation topology, equivalently the $L^1(\mu)$ topology on densities.
\end{proof}
	
	\section{Conclusion}
	
	The $L^1$/total variation topology is the natural strong topology for dominated probability measures. It reduces to the $\ell^1$ topology in the countable discrete case and to the $L^1(\mu)$ topology for densities in the general dominated case. The Jensen--Shannon divergence generates this same topology. However, entropy is not continuous in this topology in typical infinite-dimensional spaces. Therefore any ensemble topology that requires both continuous statistical mixing and continuous entropy must refine total variation: Jensen--Shannon or $L^1$ convergence alone is not sufficient.
	
\end{document}